% Pandoc template that renders paper.md in the style of a JOSS article.
% Build: see build.sh (pandoc --citeproc → xelatex). The journal itself compiles
% paper.md with its own tooling; this file only produces a local preview PDF.
\documentclass[10pt,a4paper,onecolumn]{article}
\usepackage[top=2.4cm,bottom=2.6cm,left=2.2cm,right=2.2cm]{geometry}
\usepackage{fontspec}
\setmainfont{TeX Gyre Pagella}
\setsansfont{TeX Gyre Heros}
\setmonofont[Scale=0.9]{TeX Gyre Cursor}
\usepackage{amsmath,amssymb}
\usepackage{xcolor}
\definecolor{jossblue}{RGB}{39,86,154}
\definecolor{rule}{RGB}{200,206,216}
\definecolor{muted}{RGB}{90,97,110}
\usepackage{xurl}
\usepackage[colorlinks=true,linkcolor=jossblue,citecolor=jossblue,urlcolor=jossblue,breaklinks=true]{hyperref}
\usepackage{titlesec}
\titleformat{\section}{\sffamily\bfseries\large\color{jossblue}}{}{0pt}{}
\titlespacing*{\section}{0pt}{1.6em}{0.5em}
\usepackage{enumitem}
\setlist{nosep,leftmargin=1.4em}
\usepackage{microtype}
\usepackage{fancyhdr}
\pagestyle{fancy}
\fancyhf{}
\renewcommand{\headrulewidth}{0pt}
\fancyfoot[L]{\sffamily\footnotesize\color{muted}Pulakesh
Pradhan, (2026). GeoSurvey: an installation-free, AI-assisted field
survey application with a Google Sheets backend and built-in statistical
analysis. \emph{Journal of Open Source Software}, draft.}
\fancyfoot[R]{\sffamily\footnotesize\color{muted}\thepage}
\setlength{\parindent}{0pt}
\setlength{\parskip}{0.55em}
\providecommand{\tightlist}{\setlength{\itemsep}{0pt}\setlength{\parskip}{0pt}}
% definitions for citeproc citations
\NewDocumentCommand\citeproctext{}{}
\NewDocumentCommand\citeproc{mm}{%
  \begingroup\def\citeproctext{#2}\cite{#1}\endgroup}
\makeatletter
 % allow citations to break across lines
 \let\@cite@ofmt\@firstofone
 % avoid brackets around text for \cite:
 \def\@biblabel#1{}
 \def\@cite#1#2{{#1\if@tempswa , #2\fi}}
\makeatother
\newlength{\cslhangindent}
\setlength{\cslhangindent}{1.5em}
\newlength{\csllabelwidth}
\setlength{\csllabelwidth}{3em}
\newenvironment{CSLReferences}[2] % #1 hanging-indent, #2 entry-spacing
 {\begin{list}{}{%
  \setlength{\itemindent}{0pt}
  \setlength{\leftmargin}{0pt}
  \setlength{\parsep}{0pt}
  % turn on hanging indent if param 1 is 1
  \ifodd #1
   \setlength{\leftmargin}{\cslhangindent}
   \setlength{\itemindent}{-1\cslhangindent}
  \fi
  % set entry spacing
  \setlength{\itemsep}{#2\baselineskip}}}
 {\end{list}}
\usepackage{calc}
\newcommand{\CSLBlock}[1]{\hfill\break\parbox[t]{\linewidth}{\strut\ignorespaces#1\strut}}
\newcommand{\CSLLeftMargin}[1]{\parbox[t]{\csllabelwidth}{\strut#1\strut}}
\newcommand{\CSLRightInline}[1]{\parbox[t]{\linewidth - \csllabelwidth}{\strut#1\strut}}
\newcommand{\CSLIndent}[1]{\hspace{\cslhangindent}#1}
\begin{document}

{\sffamily\bfseries\LARGE\color{jossblue}\raggedright GeoSurvey: an
installation-free, AI-assisted field survey application with a Google
Sheets backend and built-in statistical analysis\par}
\vspace{0.9em}
{\large Pulakesh
Pradhan\textsuperscript{1}~\href{https://orcid.org/0000-0003-3103-3617}{\textcolor{jossblue}{\footnotesize\sffamily ORCID 0000-0003-3103-3617}}\par}
\vspace{0.35em}
{\small\color{muted}\textsuperscript{1}\,Department of Geography,
Ravenshaw University, Cuttack, Odisha, India\par}
\vspace{0.6em}
\begin{minipage}{\linewidth}\sffamily\footnotesize\color{muted}
\textbf{DOI:} to be assigned on acceptance \quad
\textbf{Software:} \href{https://github.com/pulakeshpradhan/geosurvey}{Repository} \textbullet{} \href{https://pulakeshpradhan.github.io/geosurvey/}{Live app}\\
\textbf{Licence:} MIT \quad \textbf{Submitted:} 21 September
2026 \quad \textbf{Status:} draft for submission to the Journal of Open Source Software
\end{minipage}
\vspace{0.4em}\hrule height 0.6pt \relax\vspace{0.9em}

\section{Summary}\label{summary}

GeoSurvey is a free, open-source application for household and
socio-economic field surveys. It runs in the web browser of any phone,
tablet or laptop without installation, and stores every submission in a
Google Sheet owned by the survey team, with photographs, audio
recordings and transcripts in the team's Google Drive. A questionnaire
is designed in a visual editor, converted from a scanned paper form, or
generated from a one-line brief by a large language model; it is then
\emph{published} so that every connected phone adopts it. In the field
the app records a GPS fix and a reverse-geocoded address automatically,
stamps photographs with time and location, queues submissions while
offline, and can pre-fill answers from photographs or a recorded
interview using a multimodal model. Collected records are streamed back
from the database and can be exported as CSV, JSON, SPSS
(\texttt{.sav}), KMZ and per-record PDF. A statistics engine that runs
inside the browser produces descriptive statistics, reliability,
exploratory factor analysis, covariance- and variance-based structural
equation models, mediation and moderation tests, and a written report
that updates whenever data change.

\section{Statement of need}\label{statement-of-need}

Small research teams --- postgraduate students, local NGOs, municipal or
panchayat offices --- routinely collect household data with paper forms
and re-enter them in a spreadsheet, or with general-purpose form
builders such as Google Forms (Google LLC, 2026) that lack location
capture, offline operation, media handling and statistical export.
Specialised platforms such as Open Data Kit (Hartung et al., 2010),
KoboToolbox (KoboToolbox, 2026) and Epicollect ({Aanensen et al.}, 2009)
solve the data-collection problem well, but require either a hosted
server or an account on a third-party service, an app installation on
every device, and form authoring in a spreadsheet dialect such as
XLSForm; the analysis of the collected data happens elsewhere, typically
in SPSS, R or SmartPLS, after several manual export steps.

GeoSurvey targets the gap between these two families. It has \textbf{no
server to run and nothing to install}: the client is a static
progressive web app served from GitHub Pages, and the backend is a
single Google Apps Script bound to a Google Sheet in the team's own
account, so the data never leave infrastructure the team already
controls. Phones are connected by scanning a QR code or opening a team
link; no typing of URLs or keys is needed. Because the analysis pipeline
is part of the same application, a team can inspect reliability and
factor structure of a scale after the first day of fieldwork and adjust
the instrument, rather than discovering problems weeks later. The
application was written for socio-economic surveys in rural India ---
its transcription and auto-fill prompts handle Indian languages and
code-switched speech, and its geocoding returns village, block and
district --- but nothing in it is specific to that setting.

\section{Functionality}\label{functionality}

\textbf{Questionnaire design and management.} Sections, fields (text,
number, date, telephone, single and multiple choice, Likert items, free
text), the grouping of Likert items into \emph{constructs} for the
analysis engine, and AI hints are edited visually. \emph{Design with AI}
turns a brief or pasted text into a full instrument; a photographed
paper questionnaire can be converted the same way. Several
questionnaires can coexist in one spreadsheet, each with a dedicated
response tab; an updated questionnaire can continue in its current tab
or start a new one so that exports remain column-consistent. A
\texttt{.geosurvey} project file (schema, team link, interview language)
is saved to Drive on every publication and can be opened on any device
to reconnect it instantly.

\textbf{Field data collection.} Automatic GPS fix with accuracy and
altitude, reverse geocoding through the Google geocoder with an
OpenStreetMap Nominatim fallback (OpenStreetMap contributors, 2026),
time- and location-stamped photographs, audio recording of the
interview, an offline submission queue with background synchronisation,
and a form PDF for respondents. \emph{Analyze \& fill} sends the
photographs, recordings and any typed transcript to a multimodal model
and maps the answers to the questionnaire, returning English values for
review together with a verbatim transcript. Three engines are supported:
Gemini (Gemini Team, 2023) with the user's own key, an open-weight Gemma
model (Gemma Team, 2025) proxied by the backend so that enumerators need
no key at all, and Chrome's built-in on-device model (desktop Chrome,
photographs and text).

\textbf{Team and access model.} Submissions carry a random device
identifier; a phone can list, re-upload files for and delete only its
own records. Every administrative action --- publishing or deleting a
questionnaire, listing all records, exporting the team data set,
updating the backend --- is checked server-side against an admin token
stored as an Apps Script property. The token is held only in session
storage on the admin's device, and one central routine forgets it the
moment the backend rejects it. The backend updates itself from the
repository through the Apps Script API on a daily trigger, carrying its
configuration over, so a deployed survey keeps receiving fixes without
any manual redeployment.

\textbf{Exports and analysis.} CSV, JSON, KMZ (placemarks with
attributes for QGIS or Google Earth), SPSS \texttt{.sav} with variable
labels, value labels and measurement levels (a self-contained writer, no
dependency), and per-record PDF with photographs and Drive links. The
analysis engine is a dependency-free Web Worker
(\texttt{stats-worker.js}) implementing descriptive statistics, Welch's
\emph{t}-test, one-way ANOVA, \(\chi^2\) tests, ordinary least squares,
Cronbach's \(\alpha\) (Cronbach, 1951), Kaiser--Meyer--Olkin and
Bartlett tests (Kaiser, 1974), exploratory factor analysis by principal
components with varimax rotation, covariance-based confirmatory factor
analysis and structural models by maximum likelihood with CFI, TLI,
RMSEA and SRMR (Hu \& Bentler, 1999), composite reliability and average
variance extracted (Fornell \& Larcker, 1981), PLS-SEM with bootstrap
confidence intervals, HTMT discriminant validity (Henseler et al.,
2015), mediation, two-stage moderation and higher-order constructs (Hair
et al., 2019). Results are rendered as tables, path diagrams and a
narrative report that can be downloaded.

\section{Design}\label{design}

The client is about 5,000 lines of plain HTML, CSS and JavaScript with
no framework and no build step, which keeps it auditable by the
researchers who deploy it and lets it run from any static host. State
lives in \texttt{localStorage} and IndexedDB; a service worker caches
the shell for offline use. The backend, \texttt{backend/Code.gs},
exposes a small JSON API (\texttt{schema}, \texttt{ping},
\texttt{geocode}, \texttt{list}, submit, \texttt{attach}, \texttt{ai},
\texttt{selfUpdate} and the admin actions) and keeps a questionnaire
registry in a \emph{Config} tab. Every served file carries a version
string and the client, service worker and backend versions are kept
aligned, so an updated site invalidates stale caches on phones
automatically.

\section{Quality control}\label{quality-control}

Thirteen browser test suites (\texttt{tests/}, 165 checks) exercise the
application end to end with Playwright (Microsoft, 2026) against an
in-memory mock of the Apps Script backend: form rendering and
validation, the offline queue and synchronisation, admin gating and key
rejection, multi-questionnaire publication and sheet selection, project
files, AI design flows with and without an engine, geocoding fallback,
PDF generation, the QR team link (decoded back with jsQR), and layout
from 320 px to desktop widths.

\section{Acknowledgements}\label{acknowledgements}

The author thanks the OpenStreetMap contributors whose data provide the
geocoding fallback, and Google for the free tiers of Apps Script, Drive
and AI Studio on which the no-cost deployment relies.

\section*{References}\label{references}
\addcontentsline{toc}{section}{References}

\protect\phantomsection\label{refs}
\begin{CSLReferences}{1}{0}
\bibitem[\citeproctext]{ref-aanensen2009}
{Aanensen, D. M., Huntley, D. M., Feil, E. J., al-Own, F., \& Spratt, B.
G.} (2009). {EpiCollect}: Linking smartphones to web applications for
epidemiology, ecology and community data collection. \emph{PLoS ONE},
\emph{4}(9), e6968. \url{https://doi.org/10.1371/journal.pone.0006968}

\bibitem[\citeproctext]{ref-cronbach1951}
Cronbach, L. J. (1951). Coefficient alpha and the internal structure of
tests. \emph{Psychometrika}, \emph{16}(3), 297--334.
\url{https://doi.org/10.1007/BF02310555}

\bibitem[\citeproctext]{ref-fornell1981}
Fornell, C., \& Larcker, D. F. (1981). Evaluating structural equation
models with unobservable variables and measurement error. \emph{Journal
of Marketing Research}, \emph{18}(1), 39--50.
\url{https://doi.org/10.1177/002224378101800104}

\bibitem[\citeproctext]{ref-gemini2023}
Gemini Team. (2023). \emph{Gemini: A family of highly capable multimodal
models}. arXiv:2312.11805.
\url{https://doi.org/10.48550/arXiv.2312.11805}

\bibitem[\citeproctext]{ref-gemma3}
Gemma Team. (2025). \emph{Gemma 3 technical report}. arXiv:2503.19786.
\url{https://doi.org/10.48550/arXiv.2503.19786}

\bibitem[\citeproctext]{ref-googleforms}
Google LLC. (2026). \emph{Google {Forms}}.
\url{https://www.google.com/forms/about/}.

\bibitem[\citeproctext]{ref-hair2019}
Hair, J. F., Risher, J. J., Sarstedt, M., \& Ringle, C. M. (2019). When
to use and how to report the results of {PLS-SEM}. \emph{European
Business Review}, \emph{31}(1), 2--24.
\url{https://doi.org/10.1108/EBR-11-2018-0203}

\bibitem[\citeproctext]{ref-hartung2010}
Hartung, C., Lerer, A., Anokwa, Y., Tseng, C., Brunette, W., \&
Borriello, G. (2010). Open {Data} {Kit}: Tools to build information
services for developing regions. \emph{Proceedings of the 4th ACM/IEEE
International Conference on Information and Communication Technologies
and Development (ICTD '10)}, 18:1--18:12.
\url{https://doi.org/10.1145/2369220.2369236}

\bibitem[\citeproctext]{ref-henseler2015}
Henseler, J., Ringle, C. M., \& Sarstedt, M. (2015). A new criterion for
assessing discriminant validity in variance-based structural equation
modeling. \emph{Journal of the Academy of Marketing Science},
\emph{43}(1), 115--135. \url{https://doi.org/10.1007/s11747-014-0403-8}

\bibitem[\citeproctext]{ref-hu1999}
Hu, L., \& Bentler, P. M. (1999). Cutoff criteria for fit indexes in
covariance structure analysis: Conventional criteria versus new
alternatives. \emph{Structural Equation Modeling: A Multidisciplinary
Journal}, \emph{6}(1), 1--55.
\url{https://doi.org/10.1080/10705519909540118}

\bibitem[\citeproctext]{ref-kaiser1974}
Kaiser, H. F. (1974). An index of factorial simplicity.
\emph{Psychometrika}, \emph{39}(1), 31--36.
\url{https://doi.org/10.1007/BF02291575}

\bibitem[\citeproctext]{ref-kobotoolbox}
KoboToolbox. (2026). \emph{{KoboToolbox}: Data collection tools for
challenging environments}. \url{https://www.kobotoolbox.org/}.

\bibitem[\citeproctext]{ref-playwright}
Microsoft. (2026). \emph{Playwright: Fast and reliable end-to-end
testing for modern web apps}. \url{https://playwright.dev/}.

\bibitem[\citeproctext]{ref-osm}
OpenStreetMap contributors. (2026). \emph{Nominatim: Open-source
geocoding with {OpenStreetMap} data}. \url{https://nominatim.org/}.

\end{CSLReferences}

\end{document}
